\documentclass[../../main.tex]{subfiles}
\begin{document}


{\bf [解]}
条件(イ)(ロ)は
\begin{align}
 \begin{cases}
 f'(x)>0 \quad (x \ge 0) \\
 f(0)=a\ (a>1)  
 \end{cases} \label{eq:1}
\end{align}
である．

そこで条件(ハ)について考えよう．
点$P(t,f(t))$での接線は$\ell: y=f'(t)(x-t)+f(t)$だからこれと垂直な法線の方程式は$m: y=-\dfrac{1}{f'(t)}(x-t)+f(t)$である．
$\ell$と$x$軸の交点が$Q$，$m$と$x$軸の交点が$R$だから
\begin{align}
 Q\left(-\dfrac{f(t)}{f'(t)}+t,0\right), \qquad R(f(t)f'(t)+t,0) \label{eq:2}
\end{align}
である．ここで\cref{eq:1}より$0 \le x$で$f(x)$は単調増加で，$f(0)>0$と合わせて$f(x)>0$である．従って常に$R$の$x$座標の方が$Q$のそれよりも大きく，線分$QR$の長さ$F(t)$は
\begin{align}
F(t)=f(t)f'(t)+\dfrac{f(t)}{f'(t)} \label{eq:3}
\end{align}
である．これを(ハ)に代入して
\begin{align}
f'(t)+\dfrac{1}{f'(t)}=\dfrac{f(t)}{f'(t)} \quad\cdots\text{①} \label{eq:4}
\end{align}
である．

\begin{figure}[htb]
\begin{center}
\begin{tikzpicture}[scale=1.5, >=stealth]

  % --- 1. 座標軸 ---
  \draw[->, thick] (-1.5, 0) -- (5, 0) node[right] {$x$};
  \draw[->, thick] (0, -0.5) -- (0, 3.5) node[above] {$y$};
  \node[below left] at (0,0) {$O$};

  % パラメータ設定 (例: a = 1.5, t = 1.2)
  \def\aVal{1.5}
  \def\tVal{1.2}

  % f(t) = 0.25*t^2 + sqrt(a-1)*t + a
  \pgfmathsetmacro{\cVal}{sqrt(\aVal - 1)}
  \pgfmathsetmacro{\ft}{0.25*\tVal*\tVal + \cVal*\tVal + \aVal} % f(t)
  \pgfmathsetmacro{\fpt}{0.5*\tVal + \cVal}                    % f'(t)

  % 座標定義
  \coordinate (P) at (\tVal, \ft);
  
  % Q の x 座標: t - f(t)/f'(t)
  \pgfmathsetmacro{\xQ}{\tVal - \ft/\fpt}
  \coordinate (Q) at (\xQ, 0);

  % R の x 座標: t + f(t)*f'(t)
  \pgfmathsetmacro{\xR}{\tVal + \ft*\fpt}
  \coordinate (R) at (\xR, 0);

  % --- 2. 曲線 y = f(x) ---
  \draw[ultra thick, blue!80!black, domain=-0.2:2.2, samples=100] 
    plot (\x, {0.25*\x*\x + \cVal*\x + \aVal});
  \node[blue!80!black, right] at (2.0, 3.2) {$y = f(x)$};

  % y 切片 f(0) = a
  \fill[black] (0, \aVal) circle (1.2pt) node[left, font=\small] {$a$};

  % --- 3. 接線 \ell と 法線 m ---
  % 接线 \ell (赤色)
  \draw[thick, red!80!black, domain=\xQ-0.3:\tVal+0.6] 
    plot (\x, {\fpt*(\x - \tVal) + \ft});
  \node[red!80!black, above left] at (\xQ+0.2, 0.3) {$\ell$};

  % 法線 m (緑色)
  \draw[thick, green!50!black] ($(P)!-0.2!(R)$) -- ($(P)!1.15!(R)$);
  \node[green!50!black, above right] at (\xR-0.4, 0.4) {$m$};

  % 直角マーク (接線 \ell と 法線 m)
  \pgfmathsetmacro{\ang}{atan(\fpt)}
  \draw[thin, gray] ($(P) + ({180+\ang}:0.25)$) -- 
                    ($(P) + ({180+\ang}:0.25) + ({270+\ang}:0.25)$) -- 
                    ($(P) + ({270+\ang}:0.25)$);

  % --- 4. 線分 QR の強調 (F(t)) ---
  \draw[ultra thick, purple, <->] (\xQ, -0.2) -- (\xR, -0.2) 
    node[midway, below, font=\small] {$F(t) = |QR|$};
  \draw[dashed, gray] (Q) -- (\xQ, -0.25);
  \draw[dashed, gray] (R) -- (\xR, -0.25);

  % --- 5. 点 P, Q, R のプロットと補助線 ---
  % 点 P(t, f(t))
  \fill[black] (P) circle (1.8pt) node[above left] {$P(t, f(t))$};
  \draw[dashed, gray] (\tVal, 0) node[below, font=\small] {$t$} -- (P);

  % 点 Q, R
  \fill[red!80!black] (Q) circle (1.8pt) node[above left, font=\small] {$Q$};
  \fill[green!50!black] (R) circle (1.8pt) node[above right, font=\small] {$R$};

\end{tikzpicture}
\end{center}
\caption{曲線 $y=f(x)$ 上の点 $P$ における接線 $\ell$、法線 $m$ と交点 $Q, R$ の関係}
\end{figure}

\paragraph{(1)}

\cref{eq:4}の両辺に$f'(t)$をかけて整理すると
\begin{align}
\{f'(t)\}^2=f(t)-1 \label{eq:5}
\end{align}
だから両辺$t$で微分して，
\begin{align}
\begin{split}
2f'(t)f''(t)&=f'(t) \\
f''(t)&=\dfrac{1}{2} \quad(\because f'(t)>0) \quad\cdots\text{②} \label{eq:6}
\end{split}
\end{align}
である．

したがって$f'(t)$は単調増加である．
\cref{eq:6}の両辺積分して
\begin{align}
f'(t)=\dfrac{1}{2}t+C \quad\cdots\text{③} \label{eq:7}
\end{align}
とおける．ただし$C$は積分定数である．
\cref{eq:1}より$f'(0)=C>0\ \cdots$④である．

引き続き\cref{eq:7}の両辺を積分して\cref{eq:1}より$f(0)=a$だから
\begin{align}
f(t)=\dfrac{1}{4}t^2+Ct+a \quad\cdots\text{⑤} \label{eq:8}
\end{align}
とおける．

\cref{eq:7,eq:8}を\cref{eq:5}に代入して
\begin{align}
\begin{split}
\left(\dfrac{1}{2}t+C\right)^2&=\dfrac{1}{4}t^2+Ct+a-1 \\
\therefore
C^2&=a-1
\end{split}
\end{align}
となる．$C>0$だから$C=\sqrt{a-1}$である．従って\cref{eq:8}に代入して
\begin{align}
 f(t)=\dfrac{1}{4}t^2+\sqrt{a-1}t+a \label{eq:9}
\end{align}
である．

さて，$f(x)$は微分可能かつ連続だから$[x,x+h]$に平均値の定理を適用して
\begin{align}
f(x+h)-f(x)=f'(p)\cdot h
\end{align}
となる$p\ (x<p<x+h)$がある．ここで$f'(x)$は単調増加だったから，$C=\sqrt{a-1}$にも注意して
\begin{align}
f(x+h)-f(x)\ge f'(0)h = Ch = \sqrt{a-1}\cdot h
\end{align}
である．

以上から題意は示された．

\paragraph{(2)}

$F(t)$の最小値を求めるため，一階微分から$F(t)$の振る舞いを調べよう．
$F(t)=\dfrac{\{f(t)\}^2}{f'(t)}$の両辺を微分すると
\begin{align}
 F'(t)=\dfrac{2f(t)\{f'(t)\}^2-f''(t)\{f(t)\}^2}{\{f'(t)\}^2} = \dfrac{f(t)}{\{f'(t)\}^2}\left\{2\{f'(t)\}^2-f(t)f''(t)\right\}
\end{align}
である．$t\ge 0$の時$f(x), f'(x)>0$だからこの符号は$\{\ \}$部の符号に等しい．これを$G(t)$とおいて\cref{eq:9}を代入して
\begin{align}
\begin{split}
 G(t)
  &=2\left\{\dfrac{1}{2}t+\sqrt{a-1}\right\}^2-\left(\dfrac{1}{4}t^2+\sqrt{a-1}\,t+a\right)\cdot\dfrac{1}{2} \\
  &=2\left[\dfrac{1}{4}t^2+\sqrt{a-1}\cdot t+(a-1)\right]-\dfrac{1}{8}t^2-\dfrac{1}{2}\sqrt{a-1}\cdot t-\dfrac{1}{2}a \\
  &=\dfrac{3}{8}t^2+\dfrac{3}{2}\sqrt{a-1}\cdot t+\dfrac{3}{2}a-2
\end{split}
\end{align}
を得る．$G(t)=0$は$t$の二次式で，判別式$D$とすると
\begin{align}
 D = \dfrac{9(a-1)}{4} - \dfrac{3(3a-4)}{4} = \dfrac{3}{4} \ge 0
\end{align}
となって，これは$a$にかかわらず異なる2実数解を持つ．そこでこれを$\alpha,\beta, (\alpha<\beta)$とする．
$G(t)$の軸は$t=-2\sqrt{a-1}<0$だから，$\alpha<\beta\le0$または$$\alpha<0\le\beta$$となる．
以下この両者について考える．

\subparagraph{$1^\circ$ $\alpha<\beta\le0$つまり$f(0)\ge 0 \therefore \dfrac{4}{3}\le a$の時}

この時は$t \ge 0$で$G(t) \ge 0$となり，$F(t)$の増減表は以下のようになる．

従って$t\ge0$で$F(t)$は単調増加で求める最小値は$t=0$の時の
\begin{align}
 \min F(t)=F(0)=\dfrac{\{f(0)\}^2}{f'(0)} = \dfrac{a^2}{C} = \dfrac{a^2}{\sqrt{a-1}}
\end{align}
である．

\subparagraph{$2^\circ$ $\alpha<0\le\beta$つまり$f(0) \le 0 \therefore 1<a\le\dfrac{4}{3}$の時}

この時の$F(t)$の増減表は以下のようになる．

\begin{table}[htb]
  \centering
\caption{$F(t)$の増減表}
\begin{tabular}{c|ccccc}
$t$ & $0$ & & $\beta$ & & \\
\hline
$F'$ & & $-$ & $0$ & $+$ & \\
\hline
$F$ & & $\searrow$ & & $\nearrow$ &
\end{tabular}
\end{table}

従って$F(t)$が最小値をとるのは$t=\beta$の時である．$G(t)=2\{f'(t)\}^2-f(t)f''(t)$より$G(\beta)=0$から
\begin{align}
 2\{f'(\beta)\}^2-f(\beta)f''(\beta) = 0
\end{align}
である．\cref{eq:6}より$f''(\beta)=\dfrac{1}{2}$だから
\begin{align}
 f(\beta) = 4\{f'(\beta)\}^2 
\end{align}
であり，これと\cref{eq:5}より
\begin{align}
 \{f'(\beta)\}^2 = f(\beta)-1
\end{align}
だから，連立して解くと
\begin{align}
 f'(\beta) = \dfrac{1}{\sqrt{3}}\quad f(\beta) = \dfrac{4}{3}
\end{align}
となり，
\begin{align}
\min F(t)=F(\beta)= \dfrac{\{f(\beta)\}^2}{f'(\beta)} = \dfrac{16}{9}\sqrt{3}
\end{align}
\casesend

以上まとめて，求める最小値は
\begin{align}
\min F(t)=
\begin{cases}\displaystyle
\dfrac{16}{9}\sqrt{3} & \left(1<a\le\dfrac{4}{3}\right) \\
\dfrac{a^2}{\sqrt{a-1}} & \left(\dfrac{4}{3}\le a\right)
\end{cases}
\end{align}
である．

{\bf [別解]}

\paragraph{(2)}
\cref{eq:5}より$F(t)=\dfrac{\{f(t)\}^2}{f'(t)}$の両辺を二乗すると
\begin{align}
 F^2(t) = \dfrac{\{f(t)\}^4}{f(t)-1}
\end{align}
である．そこで$y=f(t)$とおいて
\begin{align}
 F^2(t) = \dfrac{(y+1)^4}{y} 
\end{align}
の最小値を求めに行く方が簡単に解けるだろう．




\end{document}
